\section{Passing to a symmetric algebra}\label{sec:transfer}

The final step of the construction has a general form.  It explains why
we will need vanishing in negative degrees as well as positive ones.
Throughout this section, $R$ is a finite-dimensional algebra over our
field $k$ and $D=\Hom_k(-,k)$.

We recall the standard symmetry of a trivial extension; see
\cite[Proposition~3.1]{BessenrodtHolmZimmermann2007}.

\begin{lemma}\label{tr:trivial-symmetric}
The trivial extension $R\ltimes DR$ is a finite-dimensional symmetric
algebra.  In particular, its left and right regular modules are
injective in the categories of all unital modules.
\end{lemma}

\begin{proof}
The dual bimodule structure is
\[
 (a\lambda b)(c)=\lambda(bca)
 \qquad(a,b,c\in R,\ \lambda\in DR),
\]
and the multiplication is
\[
 (a,\lambda)(b,\mu)=(ab,a\mu+\lambda b).
\]
The bimodule identities give associativity, and $(1,0)$ is a unit.
The functional $t(a,\lambda)=\lambda(1)$ gives the pairing
\[
 t\big((a,\lambda)(b,\mu)\big)=\mu(a)+\lambda(b).
\]
This pairing is symmetric and nondegenerate, and associativity of the
algebra makes it an associative form.  Thus $R\ltimes DR$ is isomorphic
to its dual as a bimodule.

For any algebra $G$ over $k$, the natural isomorphism
$\Hom_G(L,DG)\simeq\Hom_k(L,k)$ shows that $DG$ is an injective left
$G$-module; the same argument applies on the right.  This statement
holds for arbitrary unital $L$, since vector-space duality is exact.
Applying it to the symmetric algebra $G=R\ltimes DR$ proves the final
assertion.
\end{proof}

\begin{proposition}[Transfer through the trivial extension]
\label{tr:transfer}
Let $P$ be a totally acyclic complex of finite projective left
$R$-modules, and put $Z=c_0P$.  Suppose that $Z$ is nonprojective and
\begin{equation}\label{tr:two-sided-vanishing}
 H^a\Hom_R(P,Z)=0
 \qquad\text{for }a>0\text{ and for }a\leq -2.
\end{equation}
Then
\[
 A=R\ltimes DR,\qquad M=A\otimes_R Z
\]
are a finite-dimensional symmetric algebra and a finite-dimensional
nonprojective left $A$-module satisfying
\[
 \Ext_A^a(M,M)=0\qquad(a>0).
\]
\end{proposition}

\begin{proof}
Write $\nu=DR\otimes_R-$ and set $I=\nu P$.  Finite projectivity gives
natural isomorphisms
\begin{equation}\label{tr:nakayama-evaluation}
 DR\otimes_R P_j\simeq D\Hom_R(P_j,R),
 \qquad
 \Hom_R(Z,DR\otimes_R P_j)\simeq D\Hom_R(P_j,Z).
\end{equation}
Both follow first for finite free modules and then for their direct
summands.  The module $\Hom_R(P_j,R)$ is finite projective on the right,
so $I_j$ is injective on the left.  Total acyclicity of $P$ and exactness
of $D$ show that the entire complex $I$ is exact.

As a right $R$-module, $A=R\oplus DR$.  Consequently the restriction of
$A\otimes_RP$ is the exact complex $P\oplus I$.  Every term
$A\otimes_RP_j$ is finite projective over $A$, since induction preserves
a split direct-summand presentation.  Right exactness identifies
$c_0(A\otimes_RP)$ with $A\otimes_RZ=M$.  Thus the nonnegative tail of
$A\otimes_RP$ is a projective resolution of $M$.  In particular, the
tensor exactness needed here has been proved for this resolution;
flatness of $A_R$ is not an additional hypothesis.

Put $N=\nu Z$.  Again by right exactness, $N=c_0I$.  Exactness of $I$
identifies this cokernel with $\im(I_0\to I_{-1})$.  Hence
\begin{equation}\label{tr:injective-tail}
 0\longrightarrow N\longrightarrow I_{-1}
 \longrightarrow I_{-2}\longrightarrow I_{-3}\longrightarrow\cdots
\end{equation}
is an injective coresolution.  If $C^j=\Hom_R(P_j,Z)$ with differential
$\delta^j:C^j\to C^{j+1}$, applying $\Hom_R(Z,-)$ to
\eqref{tr:injective-tail} gives, in cohomological degrees $0,1,2$,
\[
 DC^{-1}\xrightarrow{D\delta^{-2}}DC^{-2}
 \xrightarrow{D\delta^{-3}}DC^{-3}.
\]
Both arrows surrounding degree one are present.  More generally, the
term in degree $a\geq0$ is $DC^{-a-1}$.  Thus, for every $a>0$,
\begin{equation}\label{tr:negative-duality}
 \Ext_R^a(Z,N)\simeq D H^{-a-1}\Hom_R(P,Z)=0.
\end{equation}
The first of these equalities uses degree $-2$, which accounts for the
endpoint in \eqref{tr:two-sided-vanishing}.

Induction--restriction adjunction on the projective resolution already
constructed now yields
\begin{align*}
 \Ext_A^a(M,M)
 &\simeq \Ext_R^a\bigl(Z,\Res_R^A M\bigr)\\
 &\simeq \Ext_R^a(Z,Z)\oplus\Ext_R^a(Z,N)=0
 \qquad(a>0).
\end{align*}
Here the first summand vanishes by
\eqref{tr:two-sided-vanishing} and Lemma~\ref{h:comparison}, and the
second by \eqref{tr:negative-duality}.

Lemma~\ref{tr:trivial-symmetric} proves that $A$ is symmetric.  The
module $M$ is finite and unital.  Finally, if $M$ were a direct summand
of $A^r$, tensoring the splitting with the quotient $A\to R$ would make
\[
 R\otimes_A M\simeq R\otimes_A(A\otimes_RZ)\simeq Z
\]
a direct summand of $R^r$.  A tensor functor preserves a split identity
composite without any flatness assumption.  This contradicts
nonprojectivity of $Z$.
\end{proof}

We will apply Proposition~\ref{tr:transfer} to a triangular algebra
$R=\Lambda$.  The rest of the proof constructs its module $Z$ together
with the complete resolution and the two vanishing ranges in
\eqref{tr:two-sided-vanishing}.
